\documentclass[pdflatex,sn-nature]{sn-jnl}

\usepackage{graphicx}
\usepackage{amsmath,amssymb}
\usepackage{booktabs}
\usepackage{xcolor}
\usepackage[T1]{fontenc}
\usepackage{lmodern}
\usepackage{textcomp}
\usepackage{lineno}

\raggedbottom
\setcitestyle{super,open={},close={},comma,sort&compress}
\graphicspath{{../results/figures/}{../results/figures/paper/}}

\begin{document}

\title[Drift and failure in brain implants]{Shared patterns of decoder drift and electrode failure suggest simple
safeguards for implanted brain--computer interfaces}

\author*[1]{\fnm{Md. Imtiaj Alam} \sur{Sajin}}\email{imtiajsajin@gmail.com}
\author[1]{\fnm{Esm E Moula} \sur{Chowdhury Abha}}\email{esmechowdhuryabha@gmail.com}
\author[1]{\fnm{Md Wahiduzzaman} \sur{Suva}}\email{wshuvo360@gmail.com}

\affil*[1]{\orgdiv{Department of Computer Science}, \orgname{American International University-Bangladesh},
\orgaddress{\city{Dhaka}, \postcode{1229}, \country{Bangladesh}}}

\abstract{Implanted brain--computer interfaces can restore communication and movement after paralysis, but their
signals change daily and decline over years. How fast decoders lose accuracy, which corrections work and how
electrodes fail are rarely compared across individuals. Here, we show shared patterns in public long-term recordings
from three monkeys and 14 people, spanning more than 2,600 sessions. A fixed decoder kept 72--77\% of its accuracy
after one night in all four individuals tested, while long-term decline varied. Elaborate corrections that need no new
training data did not beat simple rescaling of each channel. Statistics of neural change predicted failure little better than
time since calibration, whereas changes in decoder output did better. Stronger regularization protected later days at
no cost on the same day. Most electrodes that fell silent later recovered, and electrodes at the array edge lost
activity faster. These results provide baselines, recipes and a calibrated simulator for durable interfaces.}

\keywords{brain--computer interface, intracortical recording, Utah array, neural drift, decoder recalibration,
electrode failure}

\maketitle
\linenumbers

\section*{Introduction}

Intracortical brain--computer interfaces (BCIs) record from small populations of neurons and turn their activity into
commands. They have allowed people with paralysis to move robotic arms and computer cursors\cite{hochberg2006,
hochberg2012,collinger2013,pandarinath2017}, to type by attempted handwriting\cite{willett2021} and to speak through a
decoded voice\cite{willett2023,card2024,wairagkar2025}. Recent work describes daily home use over years\cite{card2026}.
For these devices to become everyday tools, their decoders must stay accurate with little effort from the user.

The recorded signals are not stable. Within a day, signals shift as electrodes move or neurons change their
firing\cite{perge2013}. Across days, the activity seen by each electrode changes\cite{chestek2011,downey2018,
bishop2014}. Over months and years, electrodes gradually lose the ability to record spikes\cite{barrese2013,
sponheim2021,bjanes2025,hahn2026}. Most systems therefore recalibrate their decoders often, which takes time away from
the user\cite{jarosiewicz2015}. Many methods aim to reduce this cost. Decoders can be trained on data from many
days\cite{sussillo2016}, updated during use\cite{jarosiewicz2015,orsborn2012,dangi2013}, or aligned to a reference day
without new labels\cite{degenhart2020,gallego2020,ma2023,karpowicz2025,wilson2025}. Monitoring scores have been
proposed to signal when recalibration is needed\cite{pun2024}.

It is still hard to judge how well these approaches work and when they are needed. Most were tested in one dataset,
over weeks rather than years, and against different baselines\cite{karpowicz2024falcon}. Monitoring scores have been
validated by their correlation with decoder error, although both grow with time since calibration\cite{pun2024}. On
the hardware side, long-term studies report falling array yield and impedance\cite{barrese2013,sponheim2021,chen2023,
hahn2026}. The dynamics of single electrodes, such as whether a silent electrode can recover, are less well described.
A shared quantitative account of how decoders and electrodes change, measured the same way across individuals, would
set realistic baselines. It would also show which corrections matter most.

Public long-term datasets now make such an account possible. They include three years of finger decoding in one
monkey\cite{temmar2025} and reaching recordings from two monkeys in another laboratory\cite{gallego2020,perich2018}.
They also include up to 20 years of array recordings from 14 BrainGate participants\cite{hahn2026}, and closed-loop
sessions with fixed decoders in two of these participants\cite{pun2024}.

Here, we analysed these datasets with one pipeline. We trained decoders on one day and tested them on later days
after a ladder of increasingly powerful corrections. A fixed decoder lost about a quarter of its accuracy overnight in
every individual tested, although the longer-term course differed between individuals. Label-free alignment of neural
activity did not beat simple per-channel renormalization. Statistics of neural change added little beyond elapsed time, while
changes in the decoder's own output did add information. Two simple practices helped: stronger regularization, and
recalibration that shrinks toward the previous decoder with a strength chosen from past sessions. Electrode failure
followed a similar process in monkeys and humans. Most silenced electrodes later recovered, and edge electrodes lost
spiking faster. Finally, a simulator calibrated to these measurements reproduced held-out years of decoder decay and
showed that drift slows as an implant ages. We release the code, the simulator and all derived results.

\section*{Results}

\subsection*{One analysis pipeline across 17 individuals}

We assembled four public datasets (Fig.~\ref{fig:overview}a and Supplementary Table~1). The decoder analyses used
monkey N, recorded over 1,242 days with Utah arrays in motor cortex during a two-finger task\cite{temmar2025,
nason2021}. They also used monkeys C and M, recorded over three and 1.5 years during reaching\cite{gallego2020,
perich2018}, and BrainGate participant T6, with 124 closed-loop cursor sessions over 3.1 years\cite{hahn2026}. The
electrode analyses used all 20 arrays from the 14 BrainGate participants (2,289 sessions) together with monkey N. The
monitoring analyses used closed-loop sessions in which participants T11 and T5 controlled a cursor with a fixed
decoder\cite{pun2024}.

\begin{figure}[tbp]
\centering
\makebox[\textwidth][c]{\includegraphics[width=170mm]{fig1_overview.pdf}}
\caption{\textbf{Data and analysis design.} \textbf{a}, Recording sessions used in this study, by individual and
analysis, plotted against years since each individual's first session. Each tick is one session. Blue, decoder
drift (monkeys N, C and M and human participant T6); green, electrode failure (20 arrays in 14 BrainGate participants);
orange, closed-loop monitoring (T11 and T5). Counts are per analysis, and some human sessions contribute to more than
one analysis. \textbf{b}, The correction ladder. A decoder trained on day $i$ is tested on a later day $j$ after
increasingly powerful corrections. Grey, no correction; light blue, corrections that need no labels (renormalizing each
channel and aligning the dominant neural subspace); dark blue, corrections that use labelled trials from day $j$
(re-learning one gain per channel and refitting the decoder while shrinking it toward the old one). A decoder trained
on day $j$ is the reference.}
\label{fig:overview}
\end{figure}

For each individual, we trained a linear decoder on one session and tested it on later sessions at gaps of 1 to 480
days. Before testing, every channel of the later session was z-scored with that session's own mean and standard
deviation. This step, which we call renormalization, needs no labels. We call the result the renormalized fixed
decoder. We compared it with a ladder of corrections (Fig.~\ref{fig:overview}b and Methods). Label-free rungs
re-centred each channel or aligned the dominant neural subspace to the training day\cite{degenhart2020}. Supervised
rungs used labelled trials from the later day. They re-learned one gain per channel, or refitted the decoder while
shrinking it toward the old one. A decoder trained on the later day served as the reference. We report retention,
the ratio of a corrected decoder's accuracy (coefficient of determination, $R^2$) to that of the reference.

Without any correction, decoders failed. When the training day's normalization was reused, median $R^2$ was below
zero at most gaps and reached $-3.0$ in monkey N (Supplementary Table~2). Renormalization is therefore the baseline
that any further correction must beat.

\subsection*{Decoders lose a quarter of their accuracy overnight}

After one night, the renormalized fixed decoder kept 73\% of the reference accuracy in monkey N (95\% confidence
interval 62--78\%) and 72\% in monkey C (62--78\%). It kept 77\% in monkey M (71--85\%) and 73\% in participant T6
(41--87\%) (Fig.~\ref{fig:decay}a,b). Overnight retention did not differ detectably between individuals (Kruskal--Wallis test,
$H = 1.98$, d.f.~$= 3$, $P = 0.58$; one value per training session, $n = 30$, 42, 14 and 24).

\begin{figure}[tbp]
\centering
\makebox[\textwidth][c]{\includegraphics[width=170mm]{fig2_decay.pdf}}
\caption{\textbf{Decoders lose a quarter of their accuracy overnight, then decline at individual rates.} \textbf{a},
Retention of the renormalized fixed decoder ($R^2$ relative to a decoder trained on the test day) against days since
training. Lines show medians and shaded bands show 95\% confidence intervals from a cluster bootstrap over training
sessions. Blue, monkey N; orange, monkey C; green, monkey M; yellow, human participant T6. The dotted line marks the
reference. \textbf{b}, Retention after one day, with 95\% confidence intervals. \textbf{c}, Days until retention falls
to one half, with 95\% confidence intervals (logarithmic axis). \textbf{d}, Retention after re-learning one gain per
channel with 300 labelled trials from the test day, shown as in \textbf{a}. Pair counts are given in Supplementary
Table~2.}
\label{fig:decay}
\end{figure}

The longer-term course was not shared. Monkey N's decoder lost half of its accuracy within 8.1 days (3.7--10.1) and
fell below zero after a year (Fig.~\ref{fig:decay}a,c). Monkey C was similar, at 9.2 days (4.6--19.5). Monkey M
declined more slowly (33.9 days; 9.5--44.0), and participant T6 more slowly still (125 days; 56--240). The reference
decoder stayed equally good throughout, with a median $R^2$ of 0.28--0.29 at every gap in monkey N. The decline therefore
reflects a growing mismatch between the decoder and the day, not poorer recordings.

How much of the loss reflects simple changes in each channel's scale? Re-learning one gain per channel (96
parameters in monkey N) recovered 14--34\% of the loss in monkey N, 28--42\% in monkey C and 19--35\% in monkey M
(Fig.~\ref{fig:decay}d; medians per gap). In participant T6, the same correction recovered 64--81\% of the loss, and
retention stayed above 0.8 at every gap. Most of the drift in the monkeys was thus more complex than a change of
channel gains, whereas in T6 channel gains captured most of it. The T6 result agrees with an earlier report that
tuning parameters on the same electrode tend to change together between days\cite{bishop2014}.

A nonlinear decoder showed the same overnight drop. On 79 matched session pairs from monkey N, a recurrent
network\cite{hochreiter1997} was more accurate than the linear decoder within a day ($R^2$ higher by 0.128; 95\%
confidence interval 0.111--0.141). It also kept slightly more of its accuracy on later days (retention higher by 0.07;
0.01--0.11), and its decline followed a similar time course (Supplementary Table~3).

\subsection*{Label-free alignment does not beat renormalization}

Several methods realign neural activity to a reference day without labels\cite{degenhart2020,gallego2020,ma2023,
karpowicz2025,wilson2025}. Their benefits were usually measured against a fixed decoder, or against a decoder whose
channel means alone were updated\cite{ma2023,wilson2025}. We asked whether they also improve on full renormalization.

Aligning the dominant neural subspace to the training day by an orthogonal Procrustes rotation\cite{schonemann1966}
clearly beat updating channel means alone (Fig.~\ref{fig:labelfree}a,b). Across training sessions, it raised median
$R^2$ by 0.053 in monkey N ($n = 80$ sessions; two-sided Wilcoxon signed-rank test, $P = 2.4\times10^{-8}$) and by 0.120
in participant T6 ($n = 122$; $P = 2.1\times10^{-16}$). It did not beat renormalization, which was better by 0.019 in
monkey N ($P = 1.1\times10^{-10}$) and by 0.011 in T6 ($P = 1.9\times10^{-7}$). At no gap in any of the four
individuals did the median aligned decoder exceed the renormalized one (Supplementary Table~2). A variant that
estimated the alignment from stable channels only\cite{degenhart2020} did no better. Most of the benefit of alignment
over mean updating therefore comes from equalizing channel variances, which renormalization already does.

\begin{figure}[tbp]
\centering
\makebox[\textwidth][c]{\includegraphics[width=170mm]{fig3_labelfree.pdf}}
\caption{\textbf{Label-free corrections and monitoring add little to renormalization and elapsed time.} \textbf{a},
Median decoding $R^2$ in monkey N after updating channel means only (grey), aligning the dominant neural subspace
by Procrustes rotation (purple) and full renormalization (blue), with 95\% cluster-bootstrap confidence intervals.
\textbf{b}, The same for human participant T6. \textbf{c}, Offline monitoring in monkey N (537 session pairs).
Absolute Spearman correlation between decoder $R^2$ and divergence of neural features, of decoder outputs, or of both
combined, before (light blue) and after (dark blue) controlling for elapsed days. The dashed line marks the correlation
of elapsed days alone. \textbf{d}, Closed-loop monitoring in the public MINDFUL data. Leave-one-day-out mean absolute
error in predicting angular error from elapsed days alone (grey), with neural divergence added (blue), or with output
divergence added (orange). T11, 15 days and 145 windows; T5, 6 days and 73 windows.}
\label{fig:labelfree}
\end{figure}

\subsection*{Neural change adds little beyond time since calibration}

A monitor that warns when a decoder is failing could make recalibration less frequent. The MINDFUL method scores
instability as the divergence between feature distributions on a test day and a reference day, computed for neural
features and for decoder outputs\cite{pun2024}. Such scores correlate with decoder error. Both quantities also grow
with time since calibration, so a correlation alone does not show that a score adds information.

We first tested this offline in monkey N (537 session pairs). Elapsed time alone correlated with decoder accuracy
(Spearman $|\rho| = 0.67$; Fig.~\ref{fig:labelfree}c). A MINDFUL-style divergence of neural features correlated more
weakly ($|\rho| = 0.34$), and its partial correlation fell to 0.16 after controlling for elapsed time. In
cross-validated prediction of $R^2$, adding the neural and output divergences to elapsed time slightly increased the
error (mean absolute error 0.090 versus 0.085; Supplementary Table~4). A broader set of 14 label-free statistics fared no better
(Supplementary Note~1). This was not because the remaining variation was noise. Day-to-day performance left over after
the time trend was reliable (split-half reliability 0.84), yet none of these statistics improved its prediction.

We then re-analysed the public MINDFUL data, in which T11 and T5 controlled a cursor in closed loop with a fixed
decoder (Fig.~\ref{fig:labelfree}d and Supplementary Table~5). In T11 (15 days), the neural divergence correlated
strongly with angular error ($r = 0.82$), consistent with the original report\cite{pun2024}. Most of this association
was shared with elapsed time, and the partial correlation was 0.37. Adding the neural divergence to elapsed days did
not improve leave-one-day-out prediction of error (mean absolute error 17.5\textdegree{} versus 17.1\textdegree{}).
Divergence of the decoder's output behaved differently. Its partial correlation with error stayed high (0.73), and
adding it reduced the prediction error to 11.9\textdegree{}. In T5, error was barely related to elapsed time over six
days ($r = 0.17$). Output divergence alone gave the best prediction (15.0\textdegree{} versus 32.9\textdegree{} for
elapsed days).

In these data, then, changes in the decoder's output carried information about performance beyond elapsed time,
whereas changes in neural statistics did not. In closed loop, a user's corrections against a failing decoder may
appear in its output.

\subsection*{Stronger regularization protects later days}

Ridge decoders are usually tuned by cross-validation within the training day. Same-day accuracy hardly changed over
five orders of magnitude of the ridge penalty $\alpha$, whereas accuracy on later days changed a lot
(Fig.~\ref{fig:regularize}a,b). In monkey N, raising $\alpha$ from 0.1, the value used in the original dataset
paper\cite{temmar2025}, to $10^4$ left same-day $R^2$ at 0.30. It raised $R^2$ at 30 days from 0.10 to 0.14, and at 120
days from 0.04 to 0.10.

\begin{figure}[tbp]
\centering
\makebox[\textwidth][c]{\includegraphics[width=170mm]{fig4_regularize.pdf}}
\caption{\textbf{Stronger regularization protects later days at no cost on the same day.} \textbf{a}, Median decoding
$R^2$ in monkey N against the ridge penalty $\alpha$, on held-out trials of the training day (lightest) and on sessions
1, 7, 30, 120 and 480 days later (progressively darker). The dotted line marks $\alpha = 0.1$, the value used in the
original dataset paper. The dashed line marks the value chosen by the 1\% rule, the largest $\alpha$ whose same-day
$R^2$ is within 1\% of the best. \textbf{b}, The same for human participant T6. \textbf{c}, Gain in $R^2$ from the 1\%
rule over $\alpha = 0.1$, on the same day and at each gap, for every individual. Points show medians and bars show 95\%
cluster-bootstrap confidence intervals. Values are listed in Supplementary Table~6.}
\label{fig:regularize}
\end{figure}

We therefore tested a simple rule that uses training-day data only: choose the largest $\alpha$ whose same-day
accuracy stays within 1\% of the best. The rule selected $\alpha = 10^4$ in monkeys N and C and in T6, and $10^3$ in
monkey M. Compared with $\alpha = 0.1$, it improved accuracy on later days in every individual
(Fig.~\ref{fig:regularize}c and Supplementary Table~6). Per training session, the median cross-day gain in $R^2$ was
0.060 in monkey N (40 of 40 sessions improved; two-sided Wilcoxon signed-rank test, $P = 1.8\times10^{-12}$) and 0.017
in monkey C (39 of 40; $P = 7.8\times10^{-11}$). It was 0.001 in monkey M (19 of 22; $P = 2.1\times10^{-4}$) and 0.021
in T6 (40 of 40; $P = 1.8\times10^{-12}$). Same-day accuracy did not fall in any individual (median changes from 0.000 to
$+0.008$).

Same-day cross-validation cannot reveal this benefit, because its curve is flat over the relevant range. This may
explain why weakly regularized decoders remain common. It is known in machine learning that validation on training
data can mislead when the data distribution shifts\cite{sugiyama2007}. For ridge regression, however, the best
penalty under shift can move in either direction\cite{patil2024}, so the direction we observed is an empirical
property of these recordings.

\subsection*{Shrinking toward the old decoder makes recalibration safe}

When labelled data from the later day are available, how few are enough? In monkey N (86 session pairs), we refitted
the decoder with ridge regression shrunk toward the old decoder, penalizing changes in both the weights and the
intercept (Methods). This estimator is established in transfer learning\cite{kuzborskij2013}, and related updates are
used for closed-loop adaptation\cite{orsborn2012,dangi2013}. With 100 labelled trials, about 2 minutes of data, it
recovered 79--89\% of the reference accuracy at every gap (Fig.~\ref{fig:recal}a,b and Supplementary Table~7). With 300
trials it matched the reference. Retraining from scratch was worse with few trials, especially at short gaps, and
re-learning channel gains alone reached a ceiling early.

\begin{figure}[tbp]
\centering
\makebox[\textwidth][c]{\includegraphics[width=170mm]{fig5_recalibrate.pdf}}
\caption{\textbf{Shrinking toward the old decoder makes recalibration efficient and safe.} \textbf{a}, Retention in
monkey N 7 days after training, against the number of labelled trials from the test day. Blue, ridge regression shrunk
toward the old decoder (weights and intercept); orange, retraining from scratch; green, re-learning one gain per
channel. Lines show medians and bands show 95\% cluster-bootstrap confidence intervals. The dotted line marks the
renormalized fixed decoder, which uses no labels. \textbf{b}, The same at 120 days. \textbf{c}, Percentage of
recalibrations (86 session pairs) that ended more than 0.01 $R^2$ below the renormalized fixed decoder, when the
shrinkage strength was chosen by cross-validation on the available trials (grey) or from past sessions (blue).}
\label{fig:recal}
\end{figure}

The shrinkage strength matters most when data are scarce. Choosing it by cross-validation on a few trials was
unreliable. With 10 trials, 10.5\% of recalibrations ended worse than no recalibration (Fig.~\ref{fig:recal}c). We
instead chose the strength from past sessions, as the value that worked best on session pairs with other training
days. This lowered harmful recalibrations to 1.2\% with 10 trials (Fisher's exact test, $P = 0.018$) and from 8.1\% to
0\% with 20 trials ($P = 0.014$), with similar median accuracy (Supplementary Table~8).

With 50 or fewer trials, re-learning only the 32 channels that the old decoder weighted most was slightly better than
a full refit ($R^2$ higher by 0.014--0.022; Supplementary Table~9). Choosing channels by label-free change statistics
was worse than a full refit.

\subsection*{Most silenced electrodes later recover}

We next asked how single electrodes fail. We used the BrainGate release, which reports threshold-crossing rates and
impedance for every electrode in every session\cite{hahn2026}, and the matching data from monkey N. An electrode was
called active when its threshold-crossing rate exceeded 2~Hz, the definition used by BrainGate\cite{hahn2026}.

Array yield and impedance declined over years in both species (Fig.~\ref{fig:failure}a,b). Impedance fell in 17 of 18
human arrays with measurements (median Spearman $\rho$ with time $-0.92$). In monkey N, median impedance fell from
302~k$\Omega$ in the first year to 171~k$\Omega$ in the fourth ($\rho = -0.88$). These trends replicate earlier
reports\cite{barrese2013,sponheim2021,chen2023,bjanes2025,hahn2026}.

\begin{figure}[tbp]
\centering
\makebox[\textwidth][c]{\includegraphics[width=170mm]{fig6_failure.pdf}}
\caption{\textbf{Electrodes fail through a shared, often reversible process.} \textbf{a}, Percentage of electrodes with
spiking activity (threshold-crossing rate of at least 2~Hz) against years since implant, for 20 human arrays (blue) and
the two arrays of monkey N (orange). \textbf{b}, Median impedance of each array relative to its first sessions
(logarithmic axis). \textbf{c}, Percentage of silenced electrodes (below 2~Hz for at least three consecutive sessions)
that later recovered (above 4~Hz for at least three consecutive sessions), for arrays with at least five silenced
electrodes. The dashed line marks the pooled human value (80\%). \textbf{d}, Difference in decline of log
threshold-crossing rate per year between edge and interior electrodes, per array with at least 10 initially active
electrodes. Negative values mean faster decline at the edge. Points show medians and bars show bootstrap 95\%
confidence intervals over electrodes. Per-array values are listed in Supplementary Table~10.}
\label{fig:failure}
\end{figure}

Loss of spiking was usually not permanent. Of 730 initially active electrodes that fell silent for at least three
consecutive sessions in the human arrays, 584 (80\%) later recovered for at least three sessions (Fig.~\ref{fig:failure}c). In
monkey N, 23 of 31 (74\%) recovered. A two-state model of switching between active and silent states gave median rates
of 0.0098 per day from active to silent and 0.0065 per day back, across human arrays (Supplementary Table~10). The
corresponding rates in monkey N were 0.0031 and 0.0008 per day.

Changes in single electrodes were often abrupt. After session-wide changes were removed, session-to-session changes in
log firing rate were heavy-tailed. Their excess kurtosis had a median of 3.6 across human arrays and was 10.2 in
monkey N.

Electrodes on the edge of the array lost spiking faster than interior electrodes (Fig.~\ref{fig:failure}d). Among 19
arrays with at least 10 initially active electrodes, the edge decline was steeper in 14 (sign test, one-sided
$P = 0.032$). Combining per-array tests across all 20 human arrays gave $P = 2.0\times10^{-8}$ (Fisher's method; 8
arrays individually $P < 0.05$). In monkey N, the effect held within each array (one-sided Mann--Whitney test, lateral
$P < 0.001$, medial $P = 0.046$). It was absent for spike-band power ($P = 0.66$ and 0.37). Spiking near the edge was
thus lost faster, while broadband activity declined evenly. This agrees with models in which micromotion strains
tissue most at the array edges\cite{forrest2025}, and with more coating damage on outer electrodes in
explants\cite{patel2023}.

\subsection*{A calibrated simulator reproduces held-out years}

Methods for stable decoding are hard to compare without a common test bed\cite{karpowicz2024falcon}. We therefore built
a feature-level simulator of chronic change (Methods). It takes a real session and returns the same session as it
might appear days to years later. It combines three processes: mixing of signals between channels, gradual turnover
of each channel's tuning, and switching of electrodes between active and silent states. The switching rates came
directly from the electrode statistics above. The mixing and turnover parameters were calibrated so that the
simulator's correction ladder matched the real ladder of monkey N during the first 700 days.

The calibrated simulator reproduced the real ladder (loss 0.074; Fig.~\ref{fig:sim}a). We then tested it on data it
had not seen: sessions after day 700, correction rungs that were not fitted, and amounts of labelled data that were not
fitted (Fig.~\ref{fig:sim}b). Its median absolute error was 0.06--0.12 in units of retention, three to four times lower
than that of a simulator without drift (0.27--0.39).

\begin{figure}[tbp]
\centering
\makebox[\textwidth][c]{\includegraphics[width=170mm]{fig7_simulator.pdf}}
\caption{\textbf{A calibrated simulator reproduces held-out years and shows that drift slows with implant age.}
\textbf{a}, Calibration on monkey N sessions before day 700. Points show real median retention and lines show the
calibrated simulator, for the renormalized fixed decoder (blue) and after re-learning channel gains (green).
\textbf{b}, Validation on sessions after day 700. Median absolute error between simulated and real retention for rungs
used in calibration, rungs not used in calibration and amounts of labelled data not used in calibration. Blue,
calibrated simulator; grey, the same simulator with drift switched off. \textbf{c}, Parameters calibrated separately
on early ($<700$ days, light blue) and late ($\geq700$ days, dark blue) sessions: session-to-session mixing, slow mixing,
initial turnover ($\times10$) and turnover time constant (years).}
\label{fig:sim}
\end{figure}

Validation also exposed a systematic miss. In the held-out years, real decoders decayed more slowly than the simulator
predicted. The real data confirmed this directly. At 120 days, retention was 0.18 (95\% confidence interval 0.03--0.31)
for decoders trained in the third year but $-0.12$ ($-0.87$ to 0.07) for decoders trained in the first year.
Calibrating the simulator on the later period (loss 0.063) gave weaker session-to-session mixing (0.77 versus 0.96)
and slower turnover (time constant 527 versus 360 days; Fig.~\ref{fig:sim}c and Supplementary Table~11). Drift thus
slows as the implant ages, and we provide parameter sets for both periods.

Finally, we asked whether training on simulated futures makes decoders last longer. It helped modestly ($R^2$ higher by
0.010; 95\% confidence interval 0.005--0.013) but less than stronger regularization alone (0.042; 0.033--0.050;
Supplementary Table~12). For linear
decoders, the simulator is therefore most useful for testing methods and recalibration policies.

\section*{Discussion}

We measured how decoders and electrodes change over years in three monkeys and 14 people with one pipeline. Four
findings recurred across individuals. A renormalized fixed decoder lost about a quarter of its accuracy overnight.
Label-free alignment did not beat renormalization. Neural change statistics added little beyond elapsed time. And
electrode silencing was usually transient, with faster loss at array edges.

These results translate into practical steps. First, renormalize each channel daily, and compare new methods against
this baseline rather than against a fixed or mean-updated decoder. Second, choose the ridge penalty as the largest
value within 1\% of the best same-day accuracy. This costs nothing on the training day and protected later days in all
four individuals. Third, when recalibrating with few labels, shrink toward the previous decoder, including its
intercept, with a strength taken from past sessions. Fourth, monitor the decoder's output, and report monitoring
statistics after controlling for time since calibration.

Why is the overnight drop so similar across species, laboratories, tasks and signal types? The similarity points to a
common source on the scale of a day. Candidates include small shifts of the array relative to the tissue and daily
changes in neural state\cite{perge2013,downey2018}. The longer-term course, in contrast, differed more than tenfold
between individuals. It may depend on the array, the implant site, the age of the implant and the task. In T6, slow
decay and strong recovery by channel gains suggest that drift was dominated by changes in channel scale. The
closed-loop task and the labels inferred from it may also have contributed, as discussed below.

The monitoring result refines rather than contradicts earlier work\cite{pun2024}. Divergence scores do track decoder
error. In our data, however, the neural part of the score added little once elapsed time was known, so a calendar
schedule was an equally good guide. The output part did add information in closed loop. We suggest that monitoring
studies report partial correlations and compare against a baseline that uses elapsed time only.

The electrode analysis adds dynamics to known yield curves\cite{barrese2013,sponheim2021,hahn2026}. Most silent
electrodes came back, which matters for decoder design. A decoder that drops channels permanently will waste
electrodes that may later carry useful signal again. Transient silencing could reflect changes in the nearby tissue
or in the recording noise floor\cite{polikov2005,kozai2015}, and our data cannot separate these causes. The edge
effect fits mechanical models\cite{forrest2025} and explant findings\cite{patel2023}. Hahn and colleagues analysed
yield, decoding signal quality and tuning stability in the same human data, and noted edge-related impedance
changes\cite{hahn2026}. We add electrode-level revival statistics, a quantitative test of edge-biased spiking loss
across arrays and species, and the dissociation between spiking and broadband power.

The simulator turns these measurements into a tool. It reproduces held-out decoder decay, its electrode switching
follows measured rates, and its parameters show that drift slows with implant age. An earlier closed-loop simulator
fitted an exponential drift model to two weeks of data\cite{wilson2025}; ours is calibrated over years and validated on
held-out time. It reproduces feature-level statistics from one monkey, not the underlying biophysics. The same code can
be calibrated to other datasets.

Our study has limitations. All decoder analyses were offline. In closed loop, users adapt to their
decoder\cite{gilja2012}, which could slow or hide decay. Human decoder drift was measured in one participant, with
labels inferred from the direction from cursor to target. Most analyses used linear decoders. A recurrent network
showed the same overnight drop, but we tested it only in monkey N. Activity thresholds and our definitions of silencing
and revival are conventions; we used BrainGate's definition so that results are comparable. Finally, the datasets
differ in task, signal type and array placement. These differences make the shared patterns more notable, but they also
limit direct comparison of long-term rates.

Long-term public datasets made this study possible. We release our code, derived results and the calibrated simulator
so that others can test methods against realistic drift and electrode failure. Our aim is to help build implanted
interfaces that need less recalibration and keep working for longer.

\section*{Methods}

\subsection*{Datasets and ethics}

All data were public and de-identified. The original studies obtained ethical approval and informed consent, as
described in the cited publications. We collected no new data.

\textit{Monkey N} (LINK; DANDI Archive 001201\cite{data_link}; ref.~\cite{temmar2025}). One rhesus macaque had two
96-channel Utah arrays in the hand area of primary motor cortex\cite{nason2021}. The release contains 312 sessions on
303 days over 1,242 days. In each session, the monkey moved two finger groups to targets in centre-out or random-target
blocks. We used the 96 released channels (64 on the medial array and 32 on the lateral array). We used spike-band
power and threshold crossings at $-4.5$ times the root-mean-square noise, in 20~ms bins\cite{nason2020}. Outputs were
the position and velocity of the two finger groups. Electrode impedance was stored on 228 days.

\textit{Monkeys C and M} (DANDI Archive 000688\cite{data_perich}; refs.~\cite{gallego2020,perich2018}). Both monkeys
had Utah arrays in motor and premotor cortex and performed centre-out and random-target reaching. Monkey C had 68
sessions from October 2013 to October 2016, and monkey M had 28 sessions from January 2014 to June 2015. The release
contains spike-sorted units. We summed the spikes of all units on each electrode in 20~ms bins, giving 96 (C) or 192 (M)
channels. Hand position and velocity were interpolated to bin centres.

\textit{BrainGate participants} (Dryad\cite{data_braingate}; ref.~\cite{hahn2026}). The yield data contain one
recording from each of 2,289 research sessions in 14 participants with 20 arrays (A1, S1--S3, T1--T3 and T5--T11). Each
file lists, for every electrode, the threshold-crossing rate at several thresholds, the impedance and the grid
position. Nineteen arrays were in the hand area of motor cortex and one was in the middle frontal gyrus. The decoding
data for T6 contain 124 closed-loop cursor sessions over 3.1 years, from one array in motor cortex. Each session
provides spike-band power in 10~ms bins, cursor and target positions, and the periods of attempted movement.

\textit{MINDFUL} (Dryad\cite{data_mindful}; ref.~\cite{pun2024}). Closed-loop cursor blocks in which T11 (15 days) and
T5 (6 days) used a fixed decoder. Each block provides neural features in 20~ms bins (spike-band power for T11 and
threshold crossings for T5), the decoder's velocity output and the instantaneous angular error of the cursor.

\subsection*{Features and renormalization}

For monkey N and T6 we used spike-band power; T6's 10~ms bins were averaged in pairs. For monkeys C and M we used
per-electrode spike counts. Each channel was z-scored with the mean and standard deviation of the decoding session's
training segment, computed without labels. We call this renormalization. It is the deployment baseline in all
analyses.

\subsection*{Decoders and accuracy}

The linear decoder was ridge regression\cite{hoerl1970} on 8 causal lags (160~ms) of the z-scored features, with an
unpenalized intercept. Unless stated otherwise, the penalty was $\alpha = 0.1$, as in ref.~\cite{temmar2025}. The
outputs were finger position and velocity for monkey N (four outputs) and cursor position and velocity for monkeys C
and M (four outputs). For T6, the output was the intended movement direction, defined as the unit vector from cursor
to target during attempted movement (two outputs), as is common in offline analyses of closed-loop data. Accuracy was
the coefficient of determination $R^2$, averaged over outputs and computed on held-out trials of the test session. In
monkey N, the first 300 of 375 trials trained the decoder and the rest were held out, as in ref.~\cite{temmar2025}. In
the other datasets, the first 80\% of trials (or movement periods) trained the decoder. The recurrent decoder was a
long short-term memory network\cite{hochreiter1997} with 256 hidden units and 20-bin input windows (400~ms). It was
trained with Adam\cite{kingma2015} for 1,500 iterations with batches of 256 on the same 300 trials.

\subsection*{Session pairs}

For each training session $i$ and target gap $g \in \{1, 7, 30, 120, 480\}$ days, we selected the later session $j$ of
the same task type whose gap was closest to $g$, within 25\%. In monkey N, 80 training sessions were drawn at random
with a fixed seed. The full monkey N set (537 pairs) also includes gaps of 2, 4, 14, 60, 240 and 900 days and was used
for the monitoring analysis. Pair counts are listed in Supplementary Table~2.

\subsection*{Correction ladder}

Let $\mathbf{x}_j(t) \in \mathbb{R}^C$ be the raw features of test day $j$, $\boldsymbol{\mu}_d$ and
$\boldsymbol{\sigma}_d$ the channel means and standard deviations of day $d$, and $f_i$ the decoder trained on day $i$.
The rungs were as follows.

\begin{itemize}
\item \textit{No correction}: $f_i\big((\mathbf{x}_j - \boldsymbol{\mu}_i)/\boldsymbol{\sigma}_i\big)$.
\item \textit{Mean updating}: $f_i\big((\mathbf{x}_j - \boldsymbol{\mu}_j)/\boldsymbol{\sigma}_i\big)$.
\item \textit{Renormalization}: $f_i(\mathbf{z}_j)$ with $\mathbf{z}_j = (\mathbf{x}_j - \boldsymbol{\mu}_j)/
\boldsymbol{\sigma}_j$.
\item \textit{Label-free alignment}: let $\mathbf{V}_i, \mathbf{V}_j \in \mathbb{R}^{C\times k}$ hold the top $k = 16$
principal axes of the z-scored training data on each day. The orthogonal matrix $\mathbf{R}$ minimizing
$\|\mathbf{V}_j\mathbf{R} - \mathbf{V}_i\|_F$ was found by Procrustes analysis\cite{schonemann1966}. Inputs were mapped
as $\tilde{\mathbf{z}} = (\mathbf{I} - \mathbf{V}_j\mathbf{V}_j^\top)\mathbf{z}_j + \mathbf{V}_i\mathbf{R}^\top
\mathbf{V}_j^\top\mathbf{z}_j$, which rotates the dominant subspace and leaves the rest unchanged.
\item \textit{Stable-channel alignment}: as above, with $\mathbf{R}$ estimated iteratively from the 60\% of channels
with the smallest loading residuals (five iterations), following ref.~\cite{degenhart2020}.
\item \textit{Channel gains}: $\tilde{z}_c = g_c z_c$ for each channel $c$, with one offset per output. Gains were
shrunk toward 1 and offsets toward the old intercept.
\item \textit{Ridge to prior}: a new decoder $(\mathbf{W}, \mathbf{b})$ fitted on the lagged features $\mathbf{H}$ of day
$j$ by minimizing
\begin{equation}
\|\mathbf{Y} - \mathbf{H}\mathbf{W} - \mathbf{1}\mathbf{b}^\top\|_F^2 + \lambda\left(\|\mathbf{W} - \mathbf{W}_i\|_F^2
+ \|\mathbf{b} - \mathbf{b}_i\|_2^2\right),
\label{eq:prior}
\end{equation}
where $(\mathbf{W}_i, \mathbf{b}_i)$ is the old decoder.
\item \textit{Reference}: a decoder trained on day $j$'s training trials.
\end{itemize}

\noindent Supervised rungs used the first $n$ labelled trials of day $j$ ($n = 300$ unless stated). Their
regularization was chosen by five-fold cross-validation over contiguous blocks of trials, using one eigendecomposition
per fold. We also computed a full input remap, $\tilde{\mathbf{z}} = (\mathbf{I} + \mathbf{D})\mathbf{z}_j + \mathbf{h}$,
fitted by L-BFGS\cite{liu1989}. Because the remap can represent any linear decoder of the same form, we used it only
as a calibration target for the simulator, not as evidence about drift structure. Retention was the ratio of a rung's
$R^2$ to the reference $R^2$ on the same test trials. The time to half retention was found by interpolating the median
retention linearly in log days.

\subsection*{Monitoring analyses}

In monkey N, for every pair we computed Gaussian Kullback--Leibler divergences between days $i$ and $j$. These used the
projections of renormalized features on day $i$'s top ten principal axes (neural divergence) and the decoder's outputs
(output divergence). We report raw Spearman correlations with $R^2$ and partial correlations given log elapsed days. We
also report time-blocked cross-validated linear prediction of $R^2$ from elapsed days, divergence, or both. The 14
additional statistics are listed in Supplementary Note~1. Split-half reliability of the residual after the time trend
used alternating pairs of trials, because centre-out trials alternate between outward and return targets, and was
corrected with the Spearman--Brown formula.

For the MINDFUL data we followed the original definitions\cite{pun2024}. Blocks were cut into non-overlapping 60~s
windows. Neural features were z-scored with reference-day statistics and projected on the reference day's top ten
principal axes. We computed Gaussian divergences of these projections and of the two-dimensional decoder output
against the first day. The median angular error in each window was taken over non-excluded trials, and windows with
fewer than 500 valid bins were skipped. We report Pearson and partial correlations with angular error given elapsed
days. We also report leave-one-day-out linear prediction of angular error from elapsed days alone or together with
each divergence. The reference day was excluded from evaluation.

\subsection*{Regularization analysis}

Ridge decoders were trained with $\alpha \in \{0.1, 10, 10^3, 10^4, 3\times10^4, 10^5, 3\times10^5, 10^6\}$ on 40
training sessions per individual (22 for monkey M) and evaluated on held-out same-day trials and on later sessions at
the five target gaps. The 1\% rule chose the largest $\alpha$ whose median same-day $R^2$ was at least 99\% of the best
median. For each training session, the cross-day gain was the mean difference in $R^2$ between the chosen $\alpha$ and
$\alpha = 0.1$ over its later sessions.

\subsection*{Recalibration with few trials}

We used 86 monkey N pairs at gaps of 1--480 days and $n \in \{10, 20, 50, 100, 300\}$ labelled trials. For ridge to prior
(equation~\ref{eq:prior}), $\lambda$ was chosen either per pair by cross-validation on the $n$ trials, or from history.
The history policy picked, from the grid $\{0.1, 1, \ldots, 10^6\}$, the value with the best median $R^2$ over pairs from
other training sessions. A recalibration was harmful when its $R^2$ was more than 0.01 below that of the renormalized
fixed decoder. For targeted recalibration, only $k$ channels' weights were re-learned while the rest stayed shrunk to
their old values. Channels were chosen by decoder importance (the norm of their weights), by label-free change (shift
in mean plus log ratio of standard deviations), by importance-weighted change, or at random.

\subsection*{Electrode failure statistics}

Statistics were computed per array with identical definitions in humans and monkey N. An electrode was active in a
session when its threshold-crossing rate at $-4.5$ times the root-mean-square noise was at least 2~Hz. Electrodes
present in fewer than 80\% of an array's sessions were excluded. An initially active electrode (active on average over
the first three sessions) was silenced when it stayed below threshold for at least three consecutive sessions. It
revived when it later exceeded 4~Hz for at least three consecutive sessions. Switching rates were fitted by maximum
likelihood to a two-state continuous-time Markov chain. The fit used all within-electrode pairs of sessions at most 400
days apart, for electrodes active in at least three sessions. For abruptness, we took changes in $\log(\text{rate} +
0.1)$ between sessions at most 7 days apart, subtracted the median change across electrodes and computed the excess
kurtosis over initially active electrodes. For impedance, values of 4,000~k$\Omega$ or more were excluded, and the trend
was the Spearman correlation of the array median with day. For the edge effect, edge electrodes were those on the
outer ring of the grid ($10\times10$ in humans and $8\times8$ in monkey N). Each electrode's decline was the slope of
$\log(\text{rate} + 0.1)$ against years. Edge and interior slopes were compared within each array with a one-sided
Mann--Whitney test, and the per-array $P$ values were combined with Fisher's method\cite{fisher1925}. In monkey N,
session-wide events (changes shared by all channels) and sessions with inflated rates (median channel rate more than
three times its median over sessions) were removed first.

\subsection*{Simulator}

The simulator transforms the z-scored features $\mathbf{Z} \in \mathbb{R}^{T\times C}$ of a real base session into those
of a simulated session $\Delta t$ days later. It applies three processes in turn (Supplementary Note~2).

\textit{Mixing.} $\mathbf{Z} \leftarrow \mathbf{Z}(\mathbf{I} + \mathbf{E})^\top$, where $\mathbf{E} = s(\Delta t)
[(1 - \beta)\mathbf{G} + \beta\mathbf{L}]$. Here $\mathbf{G}$ is a Gaussian matrix with zero diagonal and $\mathbf{L}$
is a Gaussian matrix restricted to grid neighbours, with $\beta = 0.5$. The scale combines a session-to-session part and
a slow part, $s(\Delta t)^2 = s_0^2 + [s_\infty(1 - e^{-\Delta t/\tau_s})]^2$.

\textit{Turnover.} Each channel $c$ replaces a fraction $\rho_c$ of its signal by a unit with rotated tuning. Fractions
are drawn from a beta distribution with mean $\rho(\Delta t) = \rho_0 + (\rho_\infty - \rho_0)(1 - e^{-\Delta t/
\tau_\rho})$ and concentration 5. The replacement unit is the base session's latent trajectory (top 20 principal
components) read out through a randomly rotated loading vector, plus residual noise.

\textit{Silencing.} Each channel that is active in the base session becomes silent with probability $p(\Delta t) =
h_{\text{off}}/(h_{\text{off}} + h_{\text{on}})\,[1 - e^{-(h_{\text{off}} + h_{\text{on}})\Delta t}]$. A silent channel
keeps a fraction $\kappa = 0.37$ of its signal variance, the ratio of tuning strength of silent to active channels
measured in monkey N; the rest is replaced by noise.

The switching rates $h_{\text{off}}$ and $h_{\text{on}}$ were taken from the electrode statistics of the calibration
period. The six parameters $s_0$, $s_\infty$, $\tau_s$, $\rho_0$, $\rho_\infty$ and $\tau_\rho$ were calibrated. For 12
base sessions, a decoder trained on the real day was evaluated on simulated later days with exactly the same correction
code as for real data. Regularization was fixed at the values most often chosen by cross-validation on real pairs. The
loss was the sum of squared differences between simulated and real median retention for three rungs (renormalized,
channel gains and full input remap, $n = 300$) at five gaps. We minimized it with a 24-point Sobol search\cite{sobol1967}
followed by Nelder--Mead\cite{nelder1965}, using common random numbers. Validation used base sessions after day 700
and compared simulated and real medians for all rungs and for $n \in \{20, 50, 100\}$. As a reference, we ran the same
comparison with all drift switched off.

For augmentation, ridge decoders trained on 25 sessions after day 700 (77 pairs) were augmented with eight simulated
futures per session ($\Delta t$ log-uniform between 1 and 120 days) from the simulator calibrated before day 700. They were compared with decoders trained without
augmentation, with stronger regularization, with random channel dropout and gain noise\cite{sussillo2016}, and with
data from several past sessions.

\subsection*{Statistics}

Session pairs that share a training session are not independent. Confidence intervals for medians and paired
differences therefore came from a cluster bootstrap over training sessions\cite{field2007} (2,000 resamples,
percentile intervals). Tests on decoders used one value per training session (two-sided Wilcoxon signed-rank or
Kruskal--Wallis tests). Harmful recalibrations were compared per session pair with two-sided Fisher's exact tests.
Electrode tests used one value per array (one-sided Mann--Whitney tests combined with Fisher's method, and an exact
one-sided sign test). Exact $P$ values are reported.
Analyses used Python 3.11 with NumPy, SciPy, pandas and PyTorch.

\subsection*{Use of AI tools}

An AI assistant (Claude, Anthropic) was used to write analysis code, run analyses and draft the text under the
authors' direction. The authors reviewed the code, analyses and text and take full responsibility for the content.

\section*{Data availability}

All data analysed in this study are public. LINK (monkey N) is available from the DANDI Archive at
\url{https://doi.org/10.48324/dandi.001201/0.251023.2336}. The recordings of monkeys C and M are available at
\url{https://doi.org/10.48324/dandi.000688/0.250122.1735}. The BrainGate yield and decoding data are available from
Dryad at \url{https://doi.org/10.5061/dryad.x0k6djj1h}. The MINDFUL data are available from Dryad at
\url{https://doi.org/10.5061/dryad.n2z34tn5s}. All derived results underlying the figures and tables are provided in
the code repository.

\section*{Code availability}

All code, including the simulator with its calibrated parameters, is available at
\url{https://github.com/Imtiaj-Sajin/invasive-bci} under the MIT license. A versioned archive will be deposited on
Zenodo upon publication.

\backmatter

\bmhead{Acknowledgements}
We thank the participants of the BrainGate clinical trials and the teams who recorded and released the LINK, DANDI
000688, BrainGate and MINDFUL datasets. This work received no specific funding.

\bmhead{Author contributions}
All authors contributed to the study design, analysis, interpretation and writing, and approved the final manuscript.

\bmhead{Competing interests}
The authors declare no competing interests.

\input{refs}



\end{document}
